
\documentclass{article}
\usepackage{colortbl}
\usepackage{makecell}
\usepackage{multirow}
\usepackage{supertabular}

\begin{document}

\newcounter{utterance}

\centering \large Interaction Transcript for game `cladder', experiment `full\_v1.5\_default', episode 2591 with Qwen3.5{-}27B.
\vspace{24pt}

{ \footnotesize  \setcounter{utterance}{1}
\setlength{\tabcolsep}{0pt}
\begin{supertabular}{c@{$\;$}|p{.15\linewidth}@{}p{.15\linewidth}p{.15\linewidth}p{.15\linewidth}p{.15\linewidth}p{.15\linewidth}}
    \# & \multicolumn{2}{c}{Player} && \multicolumn{2}{c}{Game Master} \\
    \hline

    \theutterance \stepcounter{utterance}  
    & & & \multicolumn{4}{p{0.6\linewidth}}{
        \cellcolor[rgb]{0.9,0.9,0.9}{
            \makecell[{{p{\linewidth}}}]{
                \texttt{\tiny{[P1$\langle$GM]}}
                \texttt{You are an expert in causal inference. The following question is not a typical commonsense query, but rather a meticulously designed question created by a professor specializing in causal inference, intended to assess the students' mastery of the course content.} \\
\texttt{Imagine a self{-}contained, hypothetical world with only the following conditions, and without any unmentioned factors or causal relationships: Gender has a direct effect on smoking and lung cancer. Smoking has a direct effect on tar deposit. Tar deposit has a direct effect on lung cancer. Gender is unobserved. For nonsmokers, the probability of high tar deposit is 77\%. For smokers, the probability of high tar deposit is 26\%. For nonsmokers and with no tar deposit, the probability of lung cancer is 71\%. For nonsmokers and with high tar deposit, the probability of lung cancer is 45\%. For smokers and with no tar deposit, the probability of lung cancer is 63\%. For smokers and with high tar deposit, the probability of lung cancer is 37\%. The overall probability of smoking is 39\%. Will smoking decrease the chance of lung cancer?} \\
\texttt{Start your answer with 'Yes' or 'No', followed by additional reasoning or evidence to support your explanation.} \\
            }
        }
    }
     \\ \\

    \theutterance \stepcounter{utterance}  
    & \multicolumn{4}{p{0.6\linewidth}}{
        \cellcolor[rgb]{0.561,0.737,0.561}{
            \makecell[{{p{\linewidth}}}]{
                \texttt{\tiny{[P1$\rangle$GM]}}
                \texttt{No} \\
\\ 
\texttt{To determine whether smoking decreases the chance of lung cancer in this hypothetical world, we must calculate the average causal effect of smoking on lung cancer, denoted as \$P(Y=1 | do(S=1)) {-} P(Y=1 | do(S=0))\$, where \$Y\$ is lung cancer and \$S\$ is smoking.} \\
\\ 
\texttt{**1. Analyze the Causal Structure and Confounding**} \\
\texttt{The problem states:} \\
\texttt{*   Gender (\$G\$) \$\textbackslash{}to\$ Smoking (\$S\$)} \\
\texttt{*   Gender (\$G\$) \$\textbackslash{}to\$ Lung Cancer (\$Y\$)} \\
\texttt{*   Smoking (\$S\$) \$\textbackslash{}to\$ Tar Deposit (\$T\$)} \\
\texttt{*   Tar Deposit (\$T\$) \$\textbackslash{}to\$ Lung Cancer (\$Y\$)} \\
\texttt{*   Gender (\$G\$) is unobserved.} \\
\\ 
\texttt{Because Gender is a common cause of both Smoking and Lung Cancer, it acts as a confounder. The observed association between smoking and lung cancer is biased. To find the true causal effect, we must use the **Backdoor Adjustment Formula** or the **g{-}formula** (Law of Total Probability over the mediator \$T\$), effectively simulating an intervention \$do(S=s)\$.} \\
\\ 
\texttt{Since \$G\$ is unobserved, we cannot condition on it directly. However, the problem provides the conditional probabilities of the outcome (\$Y\$) given the treatment (\$S\$) and the mediator (\$T\$), and the conditional probabilities of the mediator (\$T\$) given the treatment (\$S\$).} \\
\\ 
\texttt{Crucially, in the structural equation model implied here, once we intervene on \$S\$ (set \$S=s\$), the distribution of \$T\$ depends only on \$S\$ (as \$S \textbackslash{}to T\$), and the distribution of \$Y\$ depends on \$S\$ and \$T\$ (as \$S \textbackslash{}to Y\$ and \$T \textbackslash{}to Y\$). The confounding path \$S \textbackslash{}leftarrow G \textbackslash{}to Y\$ is blocked by the intervention \$do(S)\$.} \\
\\ 
\texttt{Therefore, the probability of lung cancer under an intervention \$do(S=s)\$ is calculated by summing over the possible states of the mediator \$T\$:} \\
\texttt{\$\$P(Y=1 | do(S=s)) = \textbackslash{}sum\_\{t\} P(Y=1 | S=s, T=t) \textbackslash{}cdot P(T=t | S=s)\$\$} \\
\\ 
\texttt{*Note: We do not need to know the distribution of Gender or the specific probabilities \$P(S|G)\$ or \$P(Y|G)\$ because the intervention \$do(S)\$ severs the link from \$G\$ to \$S\$, and the provided conditional probabilities \$P(T|S)\$ and \$P(Y|S,T)\$ already encapsulate the necessary mechanisms for the post{-}intervention world.*} \\
\\ 
\texttt{**2. Extract the Data**} \\
\texttt{Let \$Y=1\$ be "Lung Cancer" and \$T=1\$ be "High Tar Deposit} \\
            }
        }
    }
    & & \\ \\

    \theutterance \stepcounter{utterance}  
    & & & \multicolumn{2}{p{0.3\linewidth}}{
        \cellcolor[rgb]{0.9,0.9,0.9}{
            \makecell[{{p{\linewidth}}}]{
                \texttt{\tiny{[GM$|$GM]}}
                \texttt{No} \\
            }
        }
    }
    & & \\ \\

    \theutterance \stepcounter{utterance}  
    & & & \multicolumn{2}{p{0.3\linewidth}}{
        \cellcolor[rgb]{0.9,0.9,0.9}{
            \makecell[{{p{\linewidth}}}]{
                \texttt{\tiny{[GM$|$GM]}}
                \texttt{no} \\
            }
        }
    }
    & & \\ \\

    \theutterance \stepcounter{utterance}  
    & & & \multicolumn{2}{p{0.3\linewidth}}{
        \cellcolor[rgb]{0.9,0.9,0.9}{
            \makecell[{{p{\linewidth}}}]{
                \texttt{\tiny{[GM$|$GM]}}
                \texttt{game\_result = WIN} \\
            }
        }
    }
    & & \\ \\

\end{supertabular}
}

\end{document}
